\newif\ifuseseminar
\useseminartrue
\input{../../latex_common/header.tex.frag}

\title{Física Matemática II -- Delta de Dirac e Funções de Green}

\begin{document}
\begin{enumerate}
	\item Considere a distribuição normal
	      $p_\sigma(x)=\frac{1}{\sqrt{2\pi\sigma^2}}e^{-x^2/2\sigma^2}$, onde $\sigma>0$
	      é o desvio padrão.
	      \begin{enumerate}
		      \item Mostre que a distribuição é normalizada, i.e,
		            \begin{align}
			            \int_{-\infty}^{\infty}p_\sigma(x)\dd x = 1.
		            \end{align}
		      \item Calcule a média e o valor esperado das potências impares de $x$,
		            i.e,
		            \begin{align}
			            \left\langle x^{2n+1}\right\rangle =
			            \int_{-\infty}^{\infty}x^{2n+1}p_\sigma(x)\dd x.
		            \end{align}
		      \item Calcule a média e o valor esperado das potências pares de $x$,
		            i.e,
		            \begin{align}
			            \left\langle x^{2n}\right\rangle =
			            \int_{-\infty}^{\infty}x^{2n}p_\sigma(x)\dd x.
		            \end{align}
		      \item Usando os resultados anteriores, mostre que a integral
		            $$\int_{-\infty}^{\infty}f(x)p_\sigma(x)\dd x =
			            f(0)+\frac{\sigma^2}{2}f''(0) +
			            \frac{\sigma^4}{8}f^{(4)}(0)+\ldots.$$
		            Onde ${}'$ denota a derivada em relação a $x$.
		            Ou seja, mostre que o limite da distribuição normal quando
		            $\sigma\rightarrow 0$ é a função delta de Dirac.
	      \end{enumerate}
	\item No caso de uma dimensão, a função de Green para o operador derivada
	      segunda é proporcional a função distância $\sigma(x,y) = |x-y|$. Em três
	      dimensões, a função de Green também pode ser escrita em termos da
	      distância $$\sigma(\vec{x},\vec{y}) = \sqrt{(x_1-y_1)^2+(x_2-y_2)^2+(x_3-y_3)^2},$$
	      onde $\vec{x}=(x_1,x_2,x_3)$ e $\vec{y}=(y_1,y_2,y_3)$.
	      \begin{enumerate}
		      \item Mostre que o Laplaciano em três dimensões, i.e.,
		            $$\vec\nabla^2 = \frac{\partial^2}{\partial x_1^2}+\frac{\partial^2}{\partial x_2^2}+\frac{\partial^2}{\partial x_3^2},$$
		            atuando sobre $1/\sigma(\vec{x},\vec{y})$ é zero para todo $\vec{x}\neq\vec{y}$, i.e.,
		            $$\vec\nabla^2\frac{1}{\sigma(\vec{x},\vec{y})} = 0.$$
		      \item Regularize a função $1/\sigma(\vec{x},\vec{y})$ usando
		            o parâmetro $\varepsilon>0$, i.e,
		            $$G_\varepsilon(\vec{x},\vec{y}) = -\frac{1}{4\pi}\frac{1}{\sqrt{(x_1-y_1)^2+(x_2-y_2)^2+(x_3-y_3)^2+\varepsilon^2}}.$$
		            Mostre que a integral do Laplaciano da função regularizada é um, i.e,
		            $$\lim_{\epsilon\to0}\int_{\mathbb{R}^3}\vec\nabla^2G_\varepsilon(\vec{x},\vec{y})\dd^3\vec{x} = 1.$$
	      \end{enumerate}
\end{enumerate}
\end{document}
